% Options for packages loaded elsewhere
\PassOptionsToPackage{unicode}{hyperref}
\PassOptionsToPackage{hyphens}{url}
\PassOptionsToPackage{dvipsnames,svgnames,x11names}{xcolor}
\documentclass[
  11pt,
  a4paper,
]{article}
\usepackage{xcolor}
\usepackage[margin=1in]{geometry}
\usepackage{amsmath,amssymb}
\setcounter{secnumdepth}{-\maxdimen} % remove section numbering
\usepackage{iftex}
\ifPDFTeX
  \usepackage[T1]{fontenc}
  \usepackage[utf8]{inputenc}
  \usepackage{textcomp} % provide euro and other symbols
\else % if luatex or xetex
  \usepackage{unicode-math} % this also loads fontspec
  \defaultfontfeatures{Scale=MatchLowercase}
  \defaultfontfeatures[\rmfamily]{Ligatures=TeX,Scale=1}
\fi
\usepackage{lmodern}
\ifPDFTeX\else
  % xetex/luatex font selection
  \ifLuaTeX
    \usepackage{luaotfload}
    \directlua{luaotfload.add_fallback("mainfontfallback",{
      "NotoSerifBengali:mode=harf"
    })}
  \fi
  \setmainfont[,RawFeature={fallback=mainfontfallback}]{Latin Modern Roman}
  \ifLuaTeX
    \usepackage{luaotfload}
    \directlua{luaotfload.add_fallback("sansfontfallback",{
      "NotoSansBengali:mode=harf"
    })}
  \fi
  \setsansfont[,RawFeature={fallback=sansfontfallback}]{Latin Modern Sans}
  \ifLuaTeX
    \usepackage{luaotfload}
    \directlua{luaotfload.add_fallback("monofontfallback",{
      "NotoSansBengali:mode=harf"
    })}
  \fi
  \setmonofont[,RawFeature={fallback=monofontfallback}]{DejaVu Sans Mono}
\fi
% Use upquote if available, for straight quotes in verbatim environments
\IfFileExists{upquote.sty}{\usepackage{upquote}}{}
\IfFileExists{microtype.sty}{% use microtype if available
  \usepackage[]{microtype}
  \UseMicrotypeSet[protrusion]{basicmath} % disable protrusion for tt fonts
}{}
\makeatletter
\@ifundefined{KOMAClassName}{% if non-KOMA class
  \IfFileExists{parskip.sty}{%
    \usepackage{parskip}
  }{% else
    \setlength{\parindent}{0pt}
    \setlength{\parskip}{6pt plus 2pt minus 1pt}}
}{% if KOMA class
  \KOMAoptions{parskip=half}}
\makeatother
\setlength{\emergencystretch}{3em} % prevent overfull lines
\providecommand{\tightlist}{%
  \setlength{\itemsep}{0pt}\setlength{\parskip}{0pt}}
% paper_preamble.tex — shared LaTeX preamble for the PCMA/SI paper set.
%
% Compiled with lualatex (NOT pdflatex) because several papers contain
% Bengali script: the CLL paper realises one cognitive state into
% English, Bengali and Banglish, so the Bengali output is EVIDENCE and
% cannot be transliterated away without destroying the point being made.
%
% Font selection is NOT done here. pandoc >= 3 has native
% `mainfontfallback` / `monofontfallback` variables that emit the
% luaotfload fallback registration in the correct position relative to
% unicode-math; doing it by hand in this header lands after unicode-math
% and breaks font loading outright. See tools/build_papers.sh.

\usepackage{amsmath}
\usepackage{amssymb}
\usepackage{mathtools}   % \xleftrightarrow etc. used in cll.md §5
\usepackage{booktabs}

% ---------------------------------------------------------------------
% Unicode symbols appearing in PROSE rather than inside $...$.
%
% The papers were written as markdown and read as markdown, so they use
% real characters — "the ω-boost", "α-child", "ΔR coherence" — instead of
% LaTeX macros. Under unicode-math those codepoints are math-mode-only
% and abort the build ("\symbballowed only in math mode"). Mapping each
% to \ensuremath keeps the markdown readable on GitHub AND compilable,
% which is the whole point of keeping one source.
% ---------------------------------------------------------------------
\usepackage{newunicodechar}

% Greek, lower case
\newunicodechar{α}{\ensuremath{\alpha}}
\newunicodechar{β}{\ensuremath{\beta}}
\newunicodechar{γ}{\ensuremath{\gamma}}
\newunicodechar{δ}{\ensuremath{\delta}}
\newunicodechar{ε}{\ensuremath{\varepsilon}}
\newunicodechar{η}{\ensuremath{\eta}}
\newunicodechar{θ}{\ensuremath{\theta}}
\newunicodechar{κ}{\ensuremath{\kappa}}
\newunicodechar{λ}{\ensuremath{\lambda}}
\newunicodechar{μ}{\ensuremath{\mu}}
\newunicodechar{π}{\ensuremath{\pi}}
\newunicodechar{ρ}{\ensuremath{\rho}}
\newunicodechar{σ}{\ensuremath{\sigma}}
\newunicodechar{τ}{\ensuremath{\tau}}
\newunicodechar{φ}{\ensuremath{\varphi}}
\newunicodechar{χ}{\ensuremath{\chi}}
\newunicodechar{ψ}{\ensuremath{\psi}}
\newunicodechar{ω}{\ensuremath{\omega}}
\newunicodechar{ν}{\ensuremath{\nu}}
\newunicodechar{ξ}{\ensuremath{\xi}}
\newunicodechar{ζ}{\ensuremath{\zeta}}
\newunicodechar{ι}{\ensuremath{\iota}}

% Greek, upper case
\newunicodechar{Δ}{\ensuremath{\Delta}}
\newunicodechar{Φ}{\ensuremath{\Phi}}
\newunicodechar{Σ}{\ensuremath{\Sigma}}
\newunicodechar{Ψ}{\ensuremath{\Psi}}
\newunicodechar{Ω}{\ensuremath{\Omega}}
\newunicodechar{Γ}{\ensuremath{\Gamma}}
\newunicodechar{Λ}{\ensuremath{\Lambda}}
\newunicodechar{Θ}{\ensuremath{\Theta}}
\newunicodechar{Π}{\ensuremath{\Pi}}

% Operators and relations
\newunicodechar{×}{\ensuremath{\times}}
\newunicodechar{·}{\ensuremath{\cdot}}
\newunicodechar{±}{\ensuremath{\pm}}
\newunicodechar{−}{\ensuremath{-}}
\newunicodechar{≈}{\ensuremath{\approx}}
\newunicodechar{≤}{\ensuremath{\leq}}
\newunicodechar{≥}{\ensuremath{\geq}}
\newunicodechar{≠}{\ensuremath{\neq}}
\newunicodechar{∈}{\ensuremath{\in}}
\newunicodechar{∂}{\ensuremath{\partial}}
\newunicodechar{∇}{\ensuremath{\nabla}}
\newunicodechar{√}{\ensuremath{\sqrt{}}}
\newunicodechar{∞}{\ensuremath{\infty}}
\newunicodechar{∝}{\ensuremath{\propto}}
\newunicodechar{∑}{\ensuremath{\sum}}
\newunicodechar{∫}{\ensuremath{\int}}
\newunicodechar{ℝ}{\ensuremath{\mathbb{R}}}
\newunicodechar{ℕ}{\ensuremath{\mathbb{N}}}
\newunicodechar{‖}{\ensuremath{\|}}
\newunicodechar{⊕}{\ensuremath{\oplus}}
\newunicodechar{⊗}{\ensuremath{\otimes}}
\newunicodechar{∀}{\ensuremath{\forall}}
\newunicodechar{∃}{\ensuremath{\exists}}
\newunicodechar{¬}{\ensuremath{\neg}}
\newunicodechar{∧}{\ensuremath{\wedge}}
\newunicodechar{∨}{\ensuremath{\vee}}

% Arrows
\newunicodechar{→}{\ensuremath{\rightarrow}}
\newunicodechar{←}{\ensuremath{\leftarrow}}
\newunicodechar{↔}{\ensuremath{\leftrightarrow}}
\newunicodechar{↑}{\ensuremath{\uparrow}}
\newunicodechar{↓}{\ensuremath{\downarrow}}
\newunicodechar{⇒}{\ensuremath{\Rightarrow}}
\newunicodechar{⇐}{\ensuremath{\Leftarrow}}
\newunicodechar{⇔}{\ensuremath{\Leftrightarrow}}

% Super/subscripts written as literal characters
\newunicodechar{²}{\ensuremath{^{2}}}
\newunicodechar{³}{\ensuremath{^{3}}}
\newunicodechar{⁴}{\ensuremath{^{4}}}
\newunicodechar{⁵}{\ensuremath{^{5}}}
\newunicodechar{⁶}{\ensuremath{^{6}}}
\newunicodechar{₀}{\ensuremath{_{0}}}
\newunicodechar{₁}{\ensuremath{_{1}}}
\newunicodechar{₂}{\ensuremath{_{2}}}
\newunicodechar{₃}{\ensuremath{_{3}}}
\newunicodechar{ᵀ}{\ensuremath{^{\mathsf{T}}}}

% Status glyphs. The papers use emoji as status markers in tables
% (implemented / failed / warning). Mapped to typographic equivalents:
% a colour emoji font in lualatex is unreliable, and the MEANING is what
% carries — a reader needs to see "this one failed", not a red cross.
\newunicodechar{✅}{\ensuremath{\checkmark}}
\newunicodechar{✓}{\ensuremath{\checkmark}}
\newunicodechar{❌}{\ensuremath{\times}}
\newunicodechar{✗}{\ensuremath{\times}}
\newunicodechar{⚠}{\textbf{!}}
\newunicodechar{🧩}{\textbf{[piece]}}
\newunicodechar{📋}{\textbf{[list]}}
\newunicodechar{🟡}{\textbf{[partial]}}
\newunicodechar{🔴}{\textbf{[dead]}}
\newunicodechar{🟢}{\textbf{[live]}}
\newunicodechar{▼}{\ensuremath{\blacktriangledown}}
\newunicodechar{►}{\ensuremath{\blacktriangleright}}
\newunicodechar{□}{\ensuremath{\square}}
\newunicodechar{█}{\rule{0.6em}{0.9ex}}
\newunicodechar{•}{\textbullet}
\newunicodechar{†}{\textdagger}
\newunicodechar{‑}{-}
\newunicodechar{‐}{-}
\newunicodechar{…}{\ldots}
\newunicodechar{।}{\textbar}   % danda: rendered by the Bengali fallback in text
\newunicodechar{⁻}{\ensuremath{^{-}}}

% Zero-width formatting codepoints. These are invisible in an editor but
% are real characters to LaTeX: an emoji written as "⚠️" is TWO
% codepoints (U+26A0 U+FE0F), so declaring only the visible one leaves a
% variation selector with no definition and the build dies on a symbol
% nobody can see. Both are mapped to nothing.
\newunicodechar{^^^^fe0f}{}    % VARIATION SELECTOR-16
\newunicodechar{^^^^200d}{}    % ZERO WIDTH JOINER
\newunicodechar{^^^^200b}{}    % ZERO WIDTH SPACE

\usepackage{longtable}
\usepackage{array}

% Comparison matrices and cost ledgers are wide. Let them shrink to the
% text block rather than overflow the margin.
\usepackage{adjustbox}

% Code blocks: wrap long lines instead of running off the page. Several
% papers quote real C++ with long identifiers.
\usepackage{fvextra}
\fvset{breaklines=true,breakanywhere=true,fontsize=\small}

\setlength{\parskip}{0.5em}
\setlength{\parindent}{0pt}

\usepackage{fancyhdr}
\pagestyle{fancy}
\fancyhf{}
\fancyfoot[C]{\thepage}
\fancyhead[L]{\small\itshape Synthetic Intelligence — research lineage}
\fancyhead[R]{\small\itshape irfan@furylogic.com}
\renewcommand{\headrulewidth}{0.2pt}
\setlength{\headheight}{14pt}
\setlength{\headsep}{18pt}
\usepackage{bookmark}
\IfFileExists{xurl.sty}{\usepackage{xurl}}{} % add URL line breaks if available
\urlstyle{same}
\hypersetup{
  pdftitle={Morphological Abduction Architecture (MAA): Overcoming Topological Local Minima in Graph-Rewriting Cognition Engines},
  pdfauthor={Irfan Mahir --- Independent Researcher, Dhaka, Bangladesh --- irfan@furylogic.com},
  colorlinks=true,
  linkcolor={black},
  filecolor={Maroon},
  citecolor={Blue},
  urlcolor={[HTML]{0B5394}},
  pdfcreator={LaTeX via pandoc}}

\title{Morphological Abduction Architecture (MAA): Overcoming
Topological Local Minima in Graph-Rewriting Cognition Engines}
\author{Irfan Mahir --- Independent Researcher, Dhaka, Bangladesh ---
irfan@furylogic.com}
\date{}

\begin{document}
\maketitle

{
\hypersetup{linkcolor=}
\setcounter{tocdepth}{3}
\tableofcontents
}
\section{Morphological Abduction Architecture (MAA): Overcoming
Topological Local Minima in Graph-Rewriting Cognition
Engines}\label{morphological-abduction-architecture-maa-overcoming-topological-local-minima-in-graph-rewriting-cognition-engines}

\textbf{Author:} Irfan Mahir \emph{Independent Researcher, Dhaka,
Bangladesh} irfan@furylogic.com \textbf{Date:} July 2026
\textbf{Subjects:} Artificial Intelligence (cs.AI); Computation and
Language (cs.CL); Logic in Computer Science (cs.LO)

\begin{center}\rule{0.5\linewidth}{0.5pt}\end{center}

\subsection{Abstract}\label{abstract}

Graph-rewriting cognition engines offer a rigid, verifiable alternative
to the fluid approximations of statistical large language models.
However, when these architectures treat semantic inputs as atomic,
unbreakable nodes, they become highly susceptible to topological
stagnation. When encountering an undefined concept (an orphan node),
standard inference pipelines (deduction, induction, analogy) fail
because the search space is mathematically empty, causing the system to
infinitely loop its active cache. We propose the Morphological Abduction
Architecture (MAA), a framework that breaks the atomic constraint by
introducing Sub-Lexical Decomposition and Abductive Synthesis. By
dynamically fracturing unknown nodes and mapping their sub-components to
existing ontologies, the system is able to hypothesize structural
meaning---transforming the engine from a rigid memorizer into a dynamic
semantic synthesizer capable of state-dependent cognitive dispatching.

\begin{center}\rule{0.5\linewidth}{0.5pt}\end{center}

\subsection{1. Introduction}\label{introduction}

In discrete graph-rewriting systems, cognition is defined by the
traversal and construction of edges between known concept nodes. While
this ensures strict logical adherence, it introduces a critical
vulnerability: the ``Atomic Node Trap.''

When standard graph engines encounter a novel term---such as a
previously unseen acronym or portmanteau---the parser initializes a
singular, isolated node \(\emptyset_V\). Because this node possesses a
degree of zero (\(deg(\emptyset_V) = 0\)), there are no topological
pathways for the inference engine to traverse. Consequently, monolithic
evaluation loops (which linearly execute deduce \(\rightarrow\) induce
\(\rightarrow\) analogize pipelines) encounter a void. Lacking a
mechanism to process the unknown, the engine falls back onto its
highest-activation localized structures, resulting in pathological
repetition (e.g., infinitely looping a cached analogy about
consciousness regardless of the prompt).

To achieve general reasoning, a system must be capable of synthesizing
meaning from the unknown. The Morphological Abduction Architecture (MAA)
formalizes the mechanics of ``connecting the dots'' by physically
fragmenting orphan nodes and generating hypothesized graph structures.

\subsection{2. Component 1: Morphological
Decomposition}\label{component-1-morphological-decomposition}

To process an unknown variable, MAA dictates that a concept node must
not be treated as a fundamental particle, but as a composite structure
capable of fission.

Let the user introduce an unknown string, \(C_{unk}\). When the parser
detects that \(deg(C_{unk}) = 0\) within the active ontology \(O\), the
engine suspends standard inference and triggers a Decomposition Rule.

The string \(C_{unk}\) is fractured into a set of sub-lexical fragments
\(F = \{f_1, f_2, \dots, f_n\}\) based on linguistic boundaries,
syllables, or character clusters. The engine then projects these
fragments against the existing ontology to identify partial topological
matches:

\[ M = \{\, f_i \in F \mid \exists n \in O \text{ where } similarity(f_i, n) > \tau \,\} \]

\textbf{Case Study (The Superuser Maneuver):} If the engine encounters
the unknown node \texttt{{[}sudo{]}}, it evaluates to
\(deg([sudo]) = 0\). MAA fractures the string into \(F = \{su, do\}\).
The engine scans the ontology and successfully binds the fragments to
known concepts: \(su \rightarrow [super\_user]\)
\(do \rightarrow [execute\_action]\)

The engine has successfully extracted topological anchors from a
previously isolated node.

\subsection{3. Component 2: Abductive Synthesis
Graphing}\label{component-2-abductive-synthesis-graphing}

Once morphological anchors are established, the engine cannot use
Deduction (which requires absolute axiomatic truth) or Induction (which
requires statistical precedence). It must use \textbf{Abduction}---the
generation of the most likely explanatory hypothesis.

MAA introduces a volatile edge type to the graph schema:
\texttt{EDGE\_HYPOTHESIZED}. Unlike axiomatic edges, hypothesized edges
are temporary structural bridges used to test theoretical topologies
without corrupting the permanent ontology.

Following the extraction of match set \(M\), the Abductive Synthesis
Engine generates a Hypothesis Graph \(\Gamma_H\).

For the unknown node \texttt{{[}rdma{]}}, the engine fragments the node
into \(\{rd, ma\}\). It finds partial matches for
\(rd \rightarrow [remote / read]\) and
\(ma \rightarrow [memory / machine]\). The engine synthesizes a cluster
of theoretical connections:
\texttt{{[}rdma{]}\ -\/-EDGE\_HYPOTHESIZED-\/-\textgreater{}\ {[}remote\_memory{]}}
\texttt{{[}rdma{]}\ -\/-EDGE\_HYPOTHESIZED-\/-\textgreater{}\ {[}read\_machine{]}}

By generating these volatile edges, the node \(C_{unk}\) is no longer
orphaned. The search space is populated, allowing the cognitive loop to
resume execution, score the hypotheses against the surrounding context,
and select the highest-probability path.

\subsection{4. State-Dependent Cognitive
Dispatching}\label{state-dependent-cognitive-dispatching}

The final requirement of MAA is the abolition of the monolithic
execution loop. A static \texttt{respond()} pipeline that applies the
same sequence of rules to every input is fundamentally a memorization
behavior. True cognition requires the processing architecture to adapt
to the topology of the current graph state.

MAA replaces the static loop with a dynamic Cognitive Dispatcher
function, \(\Delta\). Before any inference is executed, \(\Delta\)
evaluates the macroscopic geometry of the active graph and routes the
execution thread to specialized sub-engines:

\begin{enumerate}
\def\labelenumi{\arabic{enumi}.}
\tightlist
\item
  \textbf{State \(\alpha\) (Topological Collision):} If the input graph
  introduces an edge that directly opposes an axiomatic invariant (e.g.,
  \texttt{machine\ -\/-has-\/-\textgreater{}\ biology}), \(\Delta\)
  routes to the \textbf{Resolution Pipeline} to isolate and reject the
  contradiction.
\item
  \textbf{State \(\beta\) (Orphan Detected):} If the input graph
  contains a node where \(deg(n) = 0\) (e.g., \texttt{{[}rdma{]}}),
  \(\Delta\) routes to the \textbf{Abductive Pipeline} (Decomposition
  and Synthesis) to map the unknown.
\item
  \textbf{State \(\gamma\) (Connected Query):} Only when the input graph
  is fully integrated into the existing ontology does \(\Delta\) route
  to the \textbf{Standard Inference Pipeline} to deduce new paths.
\end{enumerate}

\[
\Delta(\Gamma_{active}) \rightarrow
\begin{cases}
Pipeline_{Resolution} & \text{if } \exists \text{ contradiction in } \Gamma \\
Pipeline_{Abductive}  & \text{if } \exists n \in \Gamma \text{ where } deg(n) = 0 \\
Pipeline_{Standard}   & \text{otherwise}
\end{cases}
\]

\subsection{5. Conclusion}\label{conclusion}

The Morphological Abduction Architecture demonstrates that the inability
of discrete graph engines to handle unknown variables is not a
limitation of discrete mathematics, but a failure of parsing granularity
and execution routing. By allowing the engine to fracture concepts into
their morphological roots and abductively synthesize hypothesized edge
relationships, we eliminate topological local minima. Furthermore,
implementing a State-Dependent Cognitive Dispatcher ensures that the
machine actively ``reads'' its own structural state before acting,
permanently transitioning the architecture from a rigid data-retrieval
script into an adaptive, synthesizing cognition engine.

\end{document}
